% ============================================================
% MỞ ĐẦU (không đánh số chương, ngang hàng với Kết luận)
% ============================================================
\chapter*{Mở đầu}
\addcontentsline{toc}{chapter}{Mở đầu}
\label{chap:mo-dau}

\noindent\textbf{Đặt vấn đề}
\medskip

Tối ưu hóa là bài toán trung tâm của nhiều lĩnh vực khoa học kỹ thuật,
với mục tiêu tìm ra phương án tốt nhất trong một không gian khả thi
theo một tiêu chí đánh giá cho trước. Các phương pháp tối ưu cổ điển
dựa trên đạo hàm, dù có cơ sở lý thuyết chặt chẽ và hội tụ nhanh, lại
bộc lộ hai hạn chế cố hữu khi đối mặt với các bài toán thực tế: dễ mắc
kẹt tại cực trị cục bộ trên hàm mục tiêu không lồi, và đòi hỏi hàm mục
tiêu phải khả vi -- điều kiện không phải lúc nào cũng thỏa mãn, đặc
biệt với các bài toán tổ hợp như định tuyến hay lập lịch. Những hạn chế
này thúc đẩy sự phát triển của nhóm phương pháp metaheuristic dựa trên
quần thể, trong đó Thuật toán Di truyền (Genetic Algorithm -- GA) là
đại diện tiêu biểu nhất, mô phỏng nguyên lý chọn lọc tự nhiên của
Darwin để tìm kiếm nghiệm mà không cần thông tin đạo hàm hay giả định
về tính lồi của hàm mục tiêu.

\bigskip
\noindent\textbf{Mục tiêu}
\medskip

Tiểu luận đặt ra hai mục tiêu chính. Một là hệ thống hóa cơ sở lý
thuyết của GA -- từ nguồn gốc sinh học, các cách mã hóa nghiệm, đến các
thành phần và toán tử cốt lõi cấu thành một chu trình tiến hóa hoàn
chỉnh -- cùng với Lập trình Di truyền (Genetic Programming -- GP), một
mở rộng của GA cho phép tiến hóa trực tiếp cấu trúc chương trình hoặc
biểu thức thay vì chỉ tối ưu tham số số. Hai là minh họa khả năng áp
dụng của hai phương pháp này qua các thực nghiệm cụ thể, đối chiếu trực
tiếp với các mốc nghiệm tối ưu hoặc quan hệ vật lý đã biết, nhằm đánh
giá khách quan chất lượng nghiệm mà GA/GP tìm được.

\bigskip
\noindent\textbf{Phạm vi}
\medskip

Nội dung tiểu luận gồm hai chương chính và phần kết luận. Chương 1
trình bày cơ sở lý thuyết của GA và GP theo trình tự từ tổng quan bài
toán tối ưu, cơ sở khoa học, các thành phần và toán tử cốt lõi, đến quy
trình thuật toán và điều kiện dừng. Chương 2 áp dụng GA cho ba nhóm bài
toán tối ưu số -- không ràng buộc, có ràng buộc, và bài toán người du
lịch -- cùng với GP cho bài toán hồi quy symbolic trên dữ liệu thời
tiết thực tế. Phần Kết luận tổng hợp kết quả đạt được, đồng thời nhìn
nhận các hạn chế của phạm vi thực nghiệm -- như việc mỗi cấu hình tham
số chỉ chạy với một seed cố định và quy mô bài toán còn tương đối nhỏ
-- cùng các hướng phát triển khả dĩ cho những nghiên cứu tiếp theo.
